% 当前成书排版定义。

\documentclass[10pt,twoside,openany,fontset=none]{ctexbook}
\usepackage[paperwidth=185mm,paperheight=260mm,inner=18mm,outer=14mm,top=17mm,bottom=19mm,headheight=7mm,headsep=4mm,footskip=9mm]{geometry}
\usepackage{fontspec}
\usepackage{xeCJK}
\usepackage{ragged2e}
\usepackage{xcolor}
\usepackage{fancyhdr}
\usepackage{tocloft}
\usepackage{setspace}
\usepackage{xurl}
\usepackage[hidelinks,unicode,pdfencoding=auto]{hyperref}
\usepackage{bookmark}
\usepackage{etoolbox}
\usepackage{graphicx}
\usepackage{eso-pic}
\usepackage{multicol}
\usepackage{needspace}
\newfontfamily\greekfont{DejaVu Serif}

\definecolor{accent}{HTML}{43545A}
\definecolor{softgray}{HTML}{697176}
\definecolor{rulegray}{HTML}{A8ADAF}
\definecolor{footgray}{HTML}{596166}

% 字体按可用性自动选择；不随源码包分发字体文件。
% Linux 优先锁定静态 EB Garamond TTF，避免 fontconfig 将同名字体解析到 WOFF/变量字体而使 xdvipdfmx 失败。
\IfFileExists{/usr/share/fonts/truetype/ebgaramond/EBGaramond12-Regular.ttf}{%
  \setmainfont[
    Path=/usr/share/fonts/truetype/ebgaramond/,
    UprightFont=EBGaramond12-Regular.ttf,
    ItalicFont=EBGaramond12-Italic.ttf,
    BoldFont=EBGaramond12-Bold.ttf
  ]{EBGaramond12-Regular.ttf}
  \newfontfamily\latinfont[
    Path=/usr/share/fonts/truetype/ebgaramond/,
    UprightFont=EBGaramond12-Regular.ttf,
    ItalicFont=EBGaramond12-Italic.ttf,
    BoldFont=EBGaramond12-Bold.ttf
  ]{EBGaramond12-Regular.ttf}
}{%
  \IfFontExistsTF{EB Garamond}{%
    \setmainfont{EB Garamond}
    \newfontfamily\latinfont{EB Garamond}
  }{%
    \setmainfont{TeX Gyre Pagella}
    \newfontfamily\latinfont{TeX Gyre Pagella}
  }%
}

\IfFontExistsTF{Noto Serif CJK SC}{%
  \setCJKmainfont{Noto Serif CJK SC}[BoldFont=Noto Serif CJK SC Bold,ItalicFont=Noto Serif CJK SC]
  \setCJKsansfont{Noto Sans CJK SC}[BoldFont=Noto Sans CJK SC Bold]
  \newCJKfontfamily\titlecjk{Noto Serif CJK SC}[BoldFont=Noto Serif CJK SC Bold]
  \newCJKfontfamily\sanscjk{Noto Sans CJK SC}[BoldFont=Noto Sans CJK SC Bold]
}{%
  \IfFontExistsTF{Source Han Serif SC}{%
    \setCJKmainfont{Source Han Serif SC}[BoldFont=Source Han Serif SC Bold,ItalicFont=Source Han Serif SC]
    \setCJKsansfont{Source Han Sans SC}[BoldFont=Source Han Sans SC Bold]
    \newCJKfontfamily\titlecjk{Source Han Serif SC}[BoldFont=Source Han Serif SC Bold]
    \newCJKfontfamily\sanscjk{Source Han Sans SC}[BoldFont=Source Han Sans SC Bold]
  }{%
    \setCJKmainfont{FandolSong}
    \setCJKsansfont{FandolHei}
    \newCJKfontfamily\titlecjk{FandolSong}
    \newCJKfontfamily\sanscjk{FandolHei}
  }%
}
\IfFontExistsTF{AR PL Mingti2L Big5}{%
  \setCJKfallbackfamilyfont{\CJKrmdefault}[AutoFakeBold]{AR PL Mingti2L Big5}
  \setCJKfallbackfamilyfont{\CJKsfdefault}[AutoFakeBold]{AR PL Mingti2L Big5}
  \setCJKfallbackfamilyfont{titlecjk}[AutoFakeBold]{AR PL Mingti2L Big5}
  \setCJKfallbackfamilyfont{sanscjk}[AutoFakeBold]{AR PL Mingti2L Big5}
}{}
\xeCJKsetup{AutoFallBack=true}
\IfFontExistsTF{Noto Sans CJK SC}{%
  \IfFontExistsTF{Noto Sans CJK SC Medium}{%
    \newCJKfontfamily\glosscjk{Noto Sans CJK SC}[BoldFont=Noto Sans CJK SC Medium]%
  }{\newCJKfontfamily\glosscjk{Noto Sans CJK SC}[BoldFont=Noto Sans CJK SC Bold]}
}{%
  \IfFontExistsTF{Source Han Sans SC}{%
    \IfFontExistsTF{Source Han Sans SC Medium}{%
      \newCJKfontfamily\glosscjk{Source Han Sans SC}[BoldFont=Source Han Sans SC Medium]%
    }{\newCJKfontfamily\glosscjk{Source Han Sans SC}[BoldFont=Source Han Sans SC Bold]}
  }{%
    \newCJKfontfamily\glosscjk{FandolHei}
  }%
}


% 附录采用协调的无衬线字族；优先使用真实 Medium 中文字重。
\IfFontExistsTF{TeX Gyre Heros}{%
  \newfontfamily\glosslatin{TeX Gyre Heros}%
}{\IfFontExistsTF{Noto Sans}{%
  \newfontfamily\glosslatin{Noto Sans}%
}{\newcommand{\glosslatin}{\sffamily}}}

% 统一人名中的中文间隔号：固定使用拉丁字体的 U+00B7，避免 CJK 字体映射为日文中点“・”。
\newcommand{\nameperiod}{{\titlecjk ·}}

\XeTeXlinebreaklocale "zh"
\XeTeXlinebreakskip = 0pt plus 1pt
\emergencystretch=2em
% 已知输出例程会产生小于 3pt 的纵向 vbox 溢出；视觉预检确认无裁切。
% 仅抑制这一级别的提示，更大的纵向溢出仍会报告。
\vfuzz=3pt
\raggedbottom
\setlength{\parindent}{0pt}
\setlength{\parskip}{0pt}
\clubpenalty=8000
\widowpenalty=8000
\displaywidowpenalty=8000

\hypersetup{
  pdftitle={埃涅阿斯纪},
  pdfauthor={Novelist}, % 用户确认公开署名；豆瓣数字ID为252994365。
  pdfsubject={中文自由体诗译本与注释},
  pdfcreator={XeLaTeX}
}

\fancypagestyle{running}{%
  \fancyhf{}
  \renewcommand{\headrulewidth}{0.25pt}
  \renewcommand{\headrule}{\hbox to\headwidth{\color{rulegray}\leaders\hrule height \headrulewidth\hfill}}
  \fancyhead[LE]{\color{softgray}\sffamily\fontsize{7.8}{9}\selectfont\thepage}
  \fancyhead[CE]{\color{softgray}\sanscjk\fontsize{7.8}{9}\selectfont 埃涅阿斯纪}
  \fancyhead[CO]{\color{softgray}\sanscjk\fontsize{7.8}{9}\selectfont\nouppercase{\rightmark}}
  \fancyhead[RO]{\color{softgray}\sffamily\fontsize{7.8}{9}\selectfont\thepage}
}
\fancypagestyle{plain}{%
  \fancyhf{}
  \renewcommand{\headrulewidth}{0.25pt}
  \renewcommand{\headrule}{\hbox to\headwidth{\color{rulegray}\leaders\hrule height \headrulewidth\hfill}}
  \fancyhead[LE]{\color{softgray}\sffamily\fontsize{7.8}{9}\selectfont\thepage}
  \fancyhead[CE]{\color{softgray}\sanscjk\fontsize{7.8}{9}\selectfont 埃涅阿斯纪}
  \fancyhead[CO]{\color{softgray}\sanscjk\fontsize{7.8}{9}\selectfont\nouppercase{\rightmark}}
  \fancyhead[RO]{\color{softgray}\sffamily\fontsize{7.8}{9}\selectfont\thepage}
}
\fancypagestyle{chapterstart}{%
  \fancyhf{}
  \renewcommand{\headrulewidth}{0pt}
  \fancyfoot[C]{\color{softgray}\sffamily\fontsize{8}{9}\selectfont\thepage}
}
\fancypagestyle{frontplain}{%
  \fancyhf{}
  \renewcommand{\headrulewidth}{0pt}
  \fancyfoot[C]{\color{softgray}\sffamily\fontsize{8}{9}\selectfont\thepage}
}
\pagestyle{running}
\renewcommand{\contentsname}{目录}
\setcounter{tocdepth}{0}
\renewcommand{\cftchapfont}{\normalfont\titlecjk\color{accent}}
\renewcommand{\cftchappagefont}{\normalfont\latinfont\color{softgray}}
\renewcommand{\cftchapleader}{\cftdotfill{2.2}}
\setlength{\cftbeforechapskip}{7pt}
\setlength{\cftchapindent}{0pt}
\setlength{\cftchapnumwidth}{0pt}

\makeatletter
\newcommand{\verseline}[2]{%
  \par\noindent
  \begingroup
  \fontsize{10.2pt}{13.65pt}\selectfont
  \hangindent=3.55em\hangafter=1
  % 为提高扫读速度，所有页的原诗行号统一置于正文左侧。
  \llap{\makebox[2.75em][r]{\sanscjk\fontsize{7.4pt}{9pt}\selectfont\color{softgray}#1}\hspace{0.8em}}%
  #2\par
  \endgroup
  \penalty150
}

\newcommand{\bookchapter}[3]{%
  \clearpage
  \phantomsection
  \addcontentsline{toc}{chapter}{#1}
  \markboth{#1}{#1}
  \thispagestyle{chapterstart}
  \vspace*{19mm}
  \begin{center}
    {\latinfont\fontsize{10}{12}\selectfont\color{softgray}\addfontfeatures{LetterSpace=12} #2\par}
    \vspace{6mm}
    {\titlecjk\bfseries\fontsize{22}{28}\selectfont\color{accent}#1\par}
    \vspace{5mm}
  \end{center}
  \vspace{11mm}
}

\newcommand{\bookend}[1]{%
  \par\vspace{8mm}
  \begin{center}
    {\sanscjk\fontsize{8.5}{10}\selectfont\color{softgray}#1}
  \end{center}
  \par\vspace{5mm}
}

\newcounter{refsection}
\newcommand{\bookrefs}{%
  \refstepcounter{refsection}
  \phantomsection
  \pdfbookmark[1]{校注参考}{refs.\therefsection}
  \par\vspace{5mm}
  {\color{rulegray}\rule{\linewidth}{0.35pt}\par}
  \vspace{3mm}
  {\sanscjk\bfseries\fontsize{10.5}{13}\selectfont\color{accent}校注参考\par}
  \vspace{2mm}
}

\newcommand{\annotation}[1]{%
  \par\noindent
  \begingroup
  \fontsize{8.55pt}{11.5pt}\selectfont\color{softgray}
  \hangindent=1.4em\hangafter=1
  \llap{\hbox to 1.1em{\hfil\textcolor{rulegray}{•}}\hspace{0.3em}}#1\par
  \endgroup
  \vspace{0.7mm}
}
\makeatother

\makeatletter
% 字号、行距与颜色保持到脚注段落结束；不重定义 footnote，以保留链接及可选参数。
\renewcommand\@makefntext[1]{%
  \begingroup
  \setstretch{1}%
  \fontsize{8.35pt}{11.15pt}\selectfont\color{footgray}%
  \parindent 1.5em
  \noindent
  \hb@xt@1.5em{\hss\@makefnmark\hspace{0.35em}}#1\par
  \endgroup}
\makeatother
\renewcommand{\thefootnote}{\arabic{footnote}}
\setlength{\footnotesep}{5.2pt}
\setlength{\skip\footins}{8pt plus 2pt minus 2pt}
% 保持单条脚注完整；脚注与其标记必要时整体移至下一页。
\interfootnotelinepenalty=10000
% 页首脚注颜色基准必须在 footins 盒子进入输出例程后重新设置。
% 避免整条脚注因分页延迟而丢失原来的 color-special 状态。
\renewcommand{\footnoterule}{\kern-3pt\begingroup\color{rulegray}\hrule width 24mm height 0.35pt\endgroup\kern 6pt\color{footgray}}

% 当前地图、卷首图、目录与书签机制。

\graphicspath{{assets/plates/}{assets/maps/}}

\expandafter\def\csname platefile1\endcsname{book-01.pdf}
\expandafter\def\csname platefile2\endcsname{book-02.pdf}
\expandafter\def\csname platefile3\endcsname{book-03.pdf}
\expandafter\def\csname platefile4\endcsname{book-04.pdf}
\expandafter\def\csname platefile5\endcsname{book-05.pdf}
\expandafter\def\csname platefile6\endcsname{book-06.pdf}
\expandafter\def\csname platefile7\endcsname{book-07.pdf}
\expandafter\def\csname platefile8\endcsname{book-08.pdf}
\expandafter\def\csname platefile9\endcsname{book-09.pdf}
\expandafter\def\csname platefile10\endcsname{book-10.pdf}
\expandafter\def\csname platefile11\endcsname{book-11.pdf}
\expandafter\def\csname platefile12\endcsname{book-12.pdf}


% 当前读者版卷首画题名与统一图版署名。
\expandafter\def\csname platetitle1\endcsname{狄多与埃涅阿斯}
\expandafter\def\csname platetitle2\endcsname{埃涅阿斯离开特洛伊}
\expandafter\def\csname platetitle3\endcsname{埃涅阿斯与安德罗玛刻相会}
\expandafter\def\csname platetitle4\endcsname{墨丘利催促埃涅阿斯}
\expandafter\def\csname platetitle5\endcsname{埃涅阿斯与阿刻斯忒斯}
\expandafter\def\csname platetitle6\endcsname{埃涅阿斯与卡戎}
\expandafter\def\csname platetitle7\endcsname{阿玛塔与拉维尼娅在酒神女信徒之中}
\expandafter\def\csname platetitle8\endcsname{维纳斯将武器交给埃涅阿斯}
\expandafter\def\csname platetitle9\endcsname{尼苏斯杀死沃尔肯斯}
\expandafter\def\csname platetitle10\endcsname{埃涅阿斯与墨泽提乌斯、劳苏斯之战}
\expandafter\def\csname platetitle11\endcsname{卡米拉杀死奥努斯之子}
\expandafter\def\csname platetitle12\endcsname{埃涅阿斯与图尔努斯的最后决斗}
\newcommand{\platecredit}{Wenceslaus Hollar（据 Francis Cleyn 设计），蚀刻版画，约 1654 年}


\newcommand{\platefile}[1]{\csname platefile#1\endcsname}
\newcommand{\platetitle}[1]{\csname platetitle#1\endcsname}

\makeatletter
\newcommand{\manualtocchapter}[2]{%
  \addtocontents{toc}{\protect\contentsline{chapter}{#1}{\thepage}{#2}\protected@file@percent}%
}
\makeatother

\newcommand{\mappage}[3]{%
  \clearpage
  \phantomsection
  \pdfbookmark[0]{#2}{map.#3}%
  \manualtocchapter{#2}{map.#3.0}%
  \markboth{#2}{#2}%
  \thispagestyle{chapterstart}%
  \AddToShipoutPictureBG*{\AtPageLowerLeft{%
    % 地图 PDF 默认按右页定位；左页整体左移 4 mm，使装订侧仍留 18 mm。
    \ifodd\value{page}\else\hspace*{-4mm}\fi
    \includegraphics[width=\paperwidth]{#1}%
  }}%
  \null
  \clearpage
}

\newcommand{\platecontent}[3]{%
  \thispagestyle{chapterstart}
  \pdfbookmark[1]{卷首画：#3}{plate.book.#1}%
  \markboth{#3}{#3}%
  \vspace*{6mm}
  \begin{center}
    {\latinfont\fontsize{9.5}{11}\selectfont\color{softgray}\addfontfeatures{LetterSpace=11} #2\par}
    \vspace{5.5mm}
    \includegraphics[width=121mm,height=166mm,keepaspectratio]{\platefile{#1}}\par
    \vspace{5mm}
    {\titlecjk\bfseries\fontsize{11.5}{15}\selectfont\color{accent}\platetitle{#1}\par}
    \vspace{2.2mm}
    {\sanscjk\fontsize{7.1}{9.5}\selectfont\color{softgray}
      \platecredit\par}
  \end{center}
}

\renewcommand{\bookchapter}[3]{%
  \ifstrequal{#3}{1}{%
    \mappage{strabo-world.pdf}{地图 1：斯特拉波所述的世界}{strabo}%
    \mappage{aeneas-voyage.pdf}{地图 2：从特洛伊到拉提乌姆}{voyage}%
  }{}%
  \ifstrequal{#3}{7}{%
    \mappage{latium-tiber.pdf}{地图 3：拉提乌姆与台伯河}{latium}%
  }{}%
  \clearpage
  \phantomsection
  \pdfbookmark[0]{#1}{book.#3}%
  \platecontent{#3}{#2}{\platetitle{#3}}%
  \clearpage
  \phantomsection
  \pdfbookmark[1]{正文}{body.book.#3}%
  \manualtocchapter{#1}{body.book.#3.1}%
  \markboth{#1}{#1}%
  \thispagestyle{chapterstart}%
  \vspace*{19mm}
  \begin{center}
    {\latinfont\fontsize{10}{12}\selectfont\color{softgray}\addfontfeatures{LetterSpace=12} #2\par}
    \vspace{6mm}
    {\titlecjk\bfseries\fontsize{22}{28}\selectfont\color{accent}#1\par}
    \vspace{5mm}
  \end{center}
  \vspace{11mm}
}

\newcommand{\creditheading}[1]{%
  \par\Needspace{22mm}\addvspace{5mm}%
  {\glosslatin\glosscjk\bfseries\fontsize{10.5}{14}\selectfont\color{accent}#1\par}%
  \nopagebreak[4]\vspace{1.8mm}\nopagebreak[4]%
}
\newcommand{\creditline}[2]{%
  \par\addvspace{0.7mm}\noindent%
  \begin{minipage}[t]{18mm}%
    \vspace{0pt}\glosslatin\glosscjk\fontsize{8.8}{13}\selectfont\color{softgray}#1\par%
  \end{minipage}\hspace{4mm}%
  \begin{minipage}[t]{\dimexpr\linewidth-22mm\relax}%
    \vspace{0pt}\RaggedRight\glosslatin\glosscjk\fontsize{9}{13}\selectfont\color{black!90}#2\par%
  \end{minipage}\par%
}
\newcommand{\illustrationcredits}{%
  \clearpage\phantomsection%
  \pdfbookmark[0]{插图与地图来源}{back.illustrations}%
  \manualtocchapter{插图与地图来源}{back.illustrations.0}%
  \markboth{插图与地图来源}{插图与地图来源}%
  \thispagestyle{chapterstart}\vspace*{8mm}%
  {\glosslatin\glosscjk\bfseries\fontsize{21.5}{27}\selectfont\color{accent}插图与地图来源\par}%
  \vspace{4mm}%
  {\setstretch{1}\glosslatin\glosscjk\fontsize{9}{13}\selectfont\color{black!90}\RaggedRight%
  % 当前读者版插图与地图来源。
\creditheading{封面}
封面采用英国画家 J.\,M.\,W. Turner 1815 年油画《狄多建造迦太基，或迦太基帝国的兴起》（\textit{Dido building Carthage, or The Rise of the Carthaginian Empire}），现藏伦敦国家美术馆，馆藏编号 NG498。所用高清图像来自 Wikimedia Commons，属公有领域。

\creditheading{卷首版画}
十二幅卷首画均为 Wenceslaus Hollar 据 Francis Cleyn 设计刻制的《埃涅阿斯纪》蚀刻版画，约作于 1654 年；扫描图像来自多伦多大学 Wenceslaus Hollar Digital Collection，并经 Wikimedia Commons 提供，均属公有领域。图版经裁切，以灰阶呈现。

\creditline{第一卷}{《狄多与埃涅阿斯》}
\creditline{第二卷}{《埃涅阿斯离开特洛伊》}
\creditline{第三卷}{《埃涅阿斯与安德罗玛刻相会》}
\creditline{第四卷}{《墨丘利催促埃涅阿斯》}
\creditline{第五卷}{《埃涅阿斯与阿刻斯忒斯》}
\creditline{第六卷}{《埃涅阿斯与卡戎》}
\creditline{第七卷}{《阿玛塔与拉维尼娅在酒神女信徒之中》}
\creditline{第八卷}{《维纳斯将武器交给埃涅阿斯》}
\creditline{第九卷}{《尼苏斯杀死沃尔肯斯》}
\creditline{第十卷}{《埃涅阿斯与墨泽提乌斯、劳苏斯之战》}
\creditline{第十一卷}{《卡米拉杀死奥努斯之子》}
\creditline{第十二卷}{《埃涅阿斯与图尔努斯的最后决斗》}

\creditheading{地图}
\creditline{地图 1}{斯特拉波所述的世界。}
\creditline{地图 2}{从特洛伊到拉提乌姆。}
\creditline{地图 3}{拉提乌姆与台伯河。}
\par\addvspace{3mm}\noindent
航程图与拉提乌姆、台伯河图据诗中叙事绘制，路线与战事范围为示意。古代岸线数据来自 Ancient World Mapping Center（AWMC，ODbL），现代岸线来自 Natural Earth（公有领域）。斯特拉波世界图采用 Bunbury 1883 年重构图，图版出自《古代地理史》（\textit{A History of Ancient Geography}）第二卷。
}%
}

% Claude HTML 是末页设计源；以下 PDF 是云端排版流程可使用的中间渲染产物，并非交接输入前提。
\newcommand{\editioncolophon}{%
  \IfFileExists{assets/colophon.pdf}{%
    \clearpage
    \phantomsection
    \pdfbookmark[0]{版本记}{back.colophon}%
    \manualtocchapter{版本记}{back.colophon.0}%
    \thispagestyle{empty}%
    \AddToShipoutPictureBG*{\AtPageLowerLeft{%
      \includegraphics[width=\paperwidth]{assets/colophon.pdf}}}%
    \null\clearpage
  }{\PackageWarningNoLine{edition}{Colophon PDF missing; draft only}}%
}

